\section{Discussion}
\label{sec:discussion}

The aim of this work was to compare the usefulness and quality of the two AI programming aids Windsurf Cascade and Github Copilot and to relate them to quality in general. This section deals with the interpretation of the results presented in the last section.

Starting with the average speed of code and response generation for the two AI programming aids. Copilot shows significantly faster performance than Windsurf overall, with Github Copilot's responses and generation being 61.1\% faster than Windsurf's response generation. This significant difference in speed can also be attributed in part to the quality of the responses. Windsurf achieves an almost flawless quality rating across all iterations and runs, while Copilot achieves an average rating that is 31.3\% below the quality of Windsurf. This means that the Github Copilot responses in this evaluation are almost a third worse than the responses generated by Windsurf Cascade. Also noteworthy is the standard deviation mentioned above, which is 0.32 points for Windsurf and 1.35 points for Copilot, allowing conclusions to be drawn about the reliability and consistency of the responses from the two systems. In comparison, Windsurf achieves an answer rate that is almost three times more consistent than Copilot. This shows that even with different types and kinds of tasks, the procedure used by Windsurf Cascade is best equipped and provides consistent answers even for complex tasks. With Copilot, on the other hand, it seems as if the model behind it does not receive sufficiently structured information and is unable to provide consistent answers.


The reason for these results can be better understood by looking at the results of the evaluation of the individual score components. Here, Copilot shows significant problems in generating quality code, i.e. the system pays attention to best practices in Terraform programming or the use of variables and configuration files. In this category, Copilot 1.17 received an average score of 2 out of 2 possible points, meaning that just over half of the possible best practices or code quality were achieved and generated in the responses and code generated by GitHub Copilot. This can be attributed to the fact that Windsurf Cascade demonstrably assumes that the system should be ready for production after generation, while Copilot pays attention to executing the specified requirements, but does not do anything beyond what is described. This can also be attributed to the high variance in responses, as different iterations of the input prompts can be interpreted differently and no attempt is made to truly understand the entire system or the thinking behind it.

However, due to the increased complexity of the tasks generated by Windsurf, this system is also more prone to problems or potential errors. For example, in 4 out of 33 cases, it was not possible to automatically insert the generated code in Windsurf, and the code had to be entered manually into the files. With Github Copilot, on the other hand, only 2 out of 33 cases were not possible, and this behavior only occurred with Copilot for complex tasks, whereas with Windsurf it occurred in all categories. In addition, the execution of the validation and planning steps with Windsurf was not successful in one out of 33 cases. Furthermore, Copilot does not generate any code at all in two consecutive executions for the task category ''Complex'' with the prompt ''C2'' once again illustrating that Copilot quickly reaches its limits when it comes to a specific type of tasks. 

The results described imply that although the generation time of Windsurf Cascade is significantly longer, the quality of the results speaks for itself. The system achieves an almost flawless result with Terraform and is the clearly better choice for Terraform code generation under the conditions described. However, caution is advised, as this study only analyses a specific part of AI programming systems and does not deal with the agent mode of the two systems or with models other than Claude Sonnet 4, with which different systems may interact better or worse. Furthermore, the evaluation is not intended to suggest that Github Copilot is a poor tool, but rather that Windsurf achieves significantly more consistent and better results than Copilot under the given circumstances and is therefore more useful for experienced and inexperienced developers under the conditions described, as the quality and reliability of the results are higher and more meaningful.